\documentclass[a4paper, 12pt]{article}
\usepackage[utf8]{inputenc}
\usepackage[margin=0.75in]{geometry}
\usepackage{scrextend}
\usepackage{amsmath}
\DeclareMathOperator*{\argmin}{arg\,min}


\title{%
  Supplementary Note 1: \\ \vspace{2mm}
  Deme Inbreeding Spatial Coefficient Derivation \\ \vspace{2mm}
  \large Supplementary text accompanying the main paper \emph{***}}

\author{***}

\begin{document}

\maketitle

\section{The Problem}

We assume that there are $K \in \mathbb{Z}_{\geq 0}$ total demes and $N \in \mathbb{Z}_{\geq 0}$ total samples, such that $N \geq K$. Further, let $i$ and $j$ be indices for demes, where $ i, j \in \{1, \dots, K\} \text{ and } i \neq j$. Consistent with Mal\'ecot \cite{Malecot1970, Malecot1973}, we assume that each deme has a specific coefficient of inbreeding, $f$. Therefore, the coefficient of relatedness between two given demes, $r_{i,j}$, can be represented as the average inbreeding between the demes conditional on the geographic distance, $d$, between demes and a global demes migration rate, $M$. The geographic distance and global migration rate capture the rate of exchange of inbreeding between the two demes, which "weights" the expected average relatedness. As a result, $\mathbf{r}$, is a vector of length ${K\choose2}$, where each element is given by:

%
\begin{equation}
    r_{ij} = \left( \frac{f_i + f_j}{2} \right) exp\left\{ -\frac{d_{ij}}{M} \right\}
\end{equation}
%

While $r_{ii} = f_i$, these terms are excluded from the loss function as they do not provide tractable information for inter-deme relatedness.  Further, we only consider distinct demes, such that $i \neq j$. We solve for the coefficient of relatedness between two given demes, $r_{i,j}$, using a gradient descent approach and the below-specified loss function.


\section{Loss Function}
We use the individual-level measures of relatedness to identify $r_{i,j}$. Let the number of pairs between demes be given by the vector $\mathbf{P}$, such that:

%
\begin{equation}
{N\choose2} = \sum_{i=1}^{K-1} \sum_{j=i+1}^K P_{i,j} \textrm{.}
\end{equation}
%

Let $\mathbf{y}$ be a vector of length ${N\choose2}$ that contains the relatedness measures between pairs, with $p_{i,j} \in 1:P_{i,j}$ denoting the index for the pairs within a given cluster-set. Therefore, the loss function, $L$, for a given vector of parameters, $\theta$, which consists of $\{f_1, ..., f_K, M\}$, is defined as:
%
\begin{equation}
    \label{eq_loss}
    L(\theta) = {N\choose2}^{-1} \sum_{i=1}^{K-1} \sum_{j=i+1}^K \sum_{p=1}^{P_{i,j}} (y_{p_{ij}} - r_{ij})^2
\end{equation}
%

Moving forward, we drop ${N \choose 2}$ as this is an uninformative constant of proportionality.


\subsection{Partial Derivative of Coefficient of Inbreeding}
First, we find the partial derivative of \eqref{eq_loss} with respect to each of our coefficient of inbreeding parameters, $\{f_1, ..., f_K\}$, of interest. All $f$ values are essentially equivalent inside the loss function, so we can simplify this problem by working out the partial derivative for a representative $f_i$.

\noindent We expand the brackets on the right-hand-side of \eqref{eq_loss} as follows:
\begin{equation}
    \label{eq_RHS}
    (y_{p_{ij}} - r_{ij})^2 = y_{p_{ij}}^2 - 2y_{p_{ij}}r_{ij} + r_{ij}^2
\end{equation}


\noindent Then we substitute in $r_{ij}$ to obtain:
\small
\begin{equation}
    \label{eq_RHS_ex}
 y_{p_{ij}}^2 - 2y_{p_{ij}} \left(\frac{f_i + f_j}{2}\right) exp\left\{ -\frac{d_{ij}}{M} \right\} + \left(\frac{f_i^2 + 2f_if_j + f_j^2}{4}\right) exp\left\{ -2 \frac{d_{ij}}{M} \right\}
\end{equation}
\normalsize


Differentiating with respect to $f_i$ we obtain:

\begin{align}
    \label{eq_partial_f}
    \frac{\partial(y_{p_{ij}} - r_{ij})^2}{\partial f_i} &= - y_{p_{ij}} exp\left\{ -\frac{d_{ij}}{M} \right\} + \left(\frac{f_i + f_j}{2}\right) exp\left\{ -2\frac{d_{ij}}{M} \right\}
\end{align}

We then substitute this into the partial derivative of \eqref{eq_loss} with respect to $f_i$ to obtain:

\begin{equation}
    \label{eq_loss_f}
    \frac{\partial L(\theta)}{\partial f_i} =  \sum_{\substack{j=1 \\ j\neq i}}^K \sum_{p=1}^{P_{i,j}} \frac{\partial (y_{p_{ij}} - r_{ij})^2}{\partial f_i}
\end{equation}

The above equation is the partial derivative solution for the coefficient of inbreeding parameter and shows that $f_i$, which is fixed, depends on all $f_j$ and the associated pairs $P_{i,j}$ between the demes $i$ and $j$. As such, the symmetry and interdependence of ${f_i, f_j}$ in $r_{i,j}$ is apparent.



\subsection{Partial Derivative of Global Migration Rate}
Next, we find the partial derivative of \eqref{eq_loss} with respect to the global migration rate parameter of interest, $M$. As above, we can use \eqref{eq_RHS_ex} as a starting point, with the partial derivative with respect to $M$, producing a more complex result:


\small
\begin{equation}
    \label{eq_partial_M}
    \frac{\partial [(y_{p_{ij}} - r_{ij})^2]}{\partial M} = \frac{-2y_{p_{ij}}d_{ij}}{M^2} \left(\frac{f_i + f_j}{2}\right) exp\left\{ \frac{-d_{ij}}{M} \right\} +
    \frac{2d_{ij}}{M^2}\left(\frac{f_i^2 + 2f_if_j + f_j^2}{4}\right) exp\left\{ \frac{-2d_{ij}}{M}  \right\}
\end{equation}
\normalsize


\noindent When it comes to the loss function, every $i,j$ contains an element of $M$; therefore, we need to retain all terms in the sum. Again, we substitute \eqref{eq_partial_M} into the loss function to formulate our partial derivative solution for the global migration rate parameter and obtain:

\begin{equation}
    \label{eq_loss_M}
    \frac{\partial L(\theta)}{\partial M} =  \sum_{i=1}^{K-1} \sum_{j=i+1}^K \sum_{p=1}^{P_{i,j}} \frac{\partial (y_{p_{ij}} - r_{ij})^2}{\partial M}
\end{equation}


\subsection{Alternate Parameterization}
\subsubsection{Coefficient of Inbreeding Reparameterization}
Reparameterization was used to address two issues: (1) constraining inbreeding coefficients to probabilistic bounds and (2) near nonidentifiability. To enforce the probabilistic bounds of the inbreeding coefficients:  $0 \leq f_i \leq 1$, we reparameterized each $f_i$ using the \textit{logit} function, such that:
\begin{equation}
    g_i = \log\left( \frac{f_i}{1-f_i} \right)
\end{equation}

\noindent We then exploit the chain rule to reparameterize the inbreeding coefficients as follows:
\begin{align}
    \label{eq_loss_figi_reparam}
    \frac{\partial L(\theta)}{\partial g_i} =  \frac{\partial L(\theta)}{\partial f_i} \cdot \frac{\partial f_i}{\partial g_i} \\
    \frac{\partial f_i}{\partial g_i} = f_i \cdot (1-f_i)
\end{align}

\noindent Near nonidentifiability in the optimization landscape, manifested as highly similar relatedness surfaces $R_{ij}$ under different parameter combinations, was attributed to strong interdependence among parameters and the high dimensionality of $\{f_1,\dots,f_K\}$. To reduce dimensionality and impose biologic plausibility that helps to uncouple this interdependence, we adopt a centered mean with deme-specific deviations, akin to common practice in the multilevel modeling literature \cite{gelmanhill2007}. Within this new framework, let $h_i = \beta + \alpha_i$, where $\beta$ is the overall mean of the inbreeding coefficients, $\{f_1, f_2, \dots, f_K\}$ and $\alpha_i$ is the deme-specific deviation. In order to constrain $\alpha_i$ from pathologic drift, we impose centering such that: $\alpha_i = c_i - \bar{c}$, where $\bar{c} = \frac{1}{K} \sum_{k=1}^K c_k$.
Given our prior logit reparameterization, we now define $f_i = \textit{expit}(h_i) = \frac{1}{1 + e^{-h_i}}$. The parameter $f_i$ now consists of elements $\beta$ and $c$. Further exploiting the chain rule, we again reparameterize the inbreeding coefficients with respect to $\beta$ and $c$ as follows:


\begin{equation}
    \label{eq_loss_figibeta_reparam}
    \frac{\partial L(\theta)}{\partial \beta} =  \sum_{i=1}^K  \frac{\partial L(\theta)}{\partial f_i} \cdot \frac{\partial f_i}{\partial g_i} \cdot
    \frac{\partial g_i}{\partial \beta}
\end{equation}

\begin{equation}
    \frac{\partial g_i}{\partial \beta} = \frac{\partial h_i}{\partial \beta} = 1
\end{equation}




\noindent The deme-specific deviations produce a more involved partial derivative.  Recall that
\[ g_i = h_i = \beta + \alpha_i  = \beta + c_i - \frac{1}{K} \sum_{k=1}^K c_k. \]
such that $h_i$ and $g_i$ are equal due to the sequential application of \textit{logit} and \textit{expit} that result from independent reparameterizations. We then differentiate the loss function with respect to an arbitrary $c_\ell$, where $\ell \in \{1,\dots,K\}$ is defined as a separate index from $i$ and $j$, to find:



\begin{equation}
    \label{eq_loss_figic_reparam}
    \frac{\partial L(\theta)}{\partial c_\ell} =
    \sum_{i=1}^K
    \frac{\partial L(\theta)}{\partial f_i} \cdot \frac{\partial f_i}{\partial g_i} \cdot
    \frac{\partial g_i}{\partial c_\ell}
\end{equation}
\begin{equation}
    \label{eq_loss_figic_numerator_reparam}
     \frac{\partial g_i}{\partial c_\ell}  = \frac{\partial \left(\beta + c_i -\frac{1}{K} \sum_{k=1}^K c_k\right)}{\partial c_\ell}
\end{equation}
\begin{equation}
    \label{eq_loss_figic_final_reparam}
     \frac{\partial g_i}{\partial c_\ell}  = \delta_{i \ell} - \frac{1}{K}
\end{equation}

where $\delta_{i \ell}$ is a Kronecker delta defined by:

\begin{equation}
    \label{Kronecker_delta}
\delta_{i \ell} =
    \begin{cases}
            1, &         \text{if } i=\ell,\\
            0, &         \text{if } i\neq \ell.
    \end{cases}
\end{equation}





\vspace{6mm}

\noindent Our final gradient for the $\beta$ component of the deme-inbreeding coefficients is given by:

\begin{equation}
    \label{final_fib_grad}
    \frac{\partial L(\theta)}{\partial \beta}
    = \sum_{i=1}^K
      \left[
        \sum_{\substack{j=1 \\ j\neq i}}^K
        \sum_{p=1}^{P_{i,j}}
        \frac{\partial (y_{p_{ij}} - r_{ij})^2}{\partial f_i}
      \right]
      f_i (1-f_i).
\end{equation}


\vspace{8mm}

\noindent Similarly, the final gradient with respect to each deme-specific parameter $c_\ell$ is
\begin{equation}
    \label{final_fic_grad}
    \frac{\partial L(\theta)}{\partial c_\ell}
    = \sum_{i=1}^K
      \left[
        \sum_{\substack{j=1 \\ j\neq i}}^K
        \sum_{p=1}^{P_{i,j}}
        \frac{\partial (y_{p_{ij}} - r_{ij})^2}{\partial f_i}
      \right]
      f_i (1-f_i)\,
      \left(\delta_{i\ell} - \frac{1}{K}\right).
\end{equation}

\vspace{8mm}

\subsubsection{Global Migration Rate Reparameterization}
For biologic plausibility, we ensure that the global migration parameter is positive, $M \geq 0$, using a log transformation and use the chain rule for reparameterization:
\begin{align}
    \label{eq_loss_M_reparam}
    m = log(M) \\
    \frac{\partial L(\theta)}{\partial m} =  \frac{\partial L(\theta)}{\partial M} \cdot \frac{\partial M}{\partial m} \\
    \frac{\partial M}{\partial m} = \frac{\partial}{\partial m} e^m = e^m
\end{align}


\vspace{6mm}

Our final gradient for the global migration rate parameter is given by:
\begin{equation}
    \label{final_M_gradient}
    \frac{\partial L(\theta)}{\partial m}
    = \left[
        \sum_{i=1}^{K-1} \sum_{j=i+1}^K \sum_{p=1}^{P_{i,j}}
        \frac{\partial (y_{p_{ij}} - r_{ij})^2}{\partial M}
      \right] e^m.
\end{equation}

\vspace{8mm}

\section{Stochastic Optimization}
\subsection{ Gradient Descent}

In order to solve for the coefficients of relatedness, $r_{i,j}$, we find the minimum of the loss function (Eq. \ref{eq_loss}) using a gradient descent approach:

\begin{align*}
\argmin_{\theta} L(\theta) \\
\theta^{t+1} = \theta^t - \eta \nabla_{\theta} L(\theta_t)
\end{align*}

\noindent where minimum is identified iteratively with $t$ steps, learning rate $\eta$, and the derived gradient from the loss function considered for each component,  $\theta = \{f_1, \dots, f_K, M\}$.

\subsection{Adaptive Moment Estimation Extension}
To further mitigate the issue of near-nonidentifiability we extend the standard gradient descent algorithm with the Adaptive Moment Estimation (ADAM) algorithm \cite{kingma2014adam}, which provides several useful features for our use-case: (1) enables an adaptive learning rate, which provides flexibility between $f_i$ and $M$, such that perturbations in each parameter will have more balanced effects; (2) mitigates collinearity between $f_i$ parameters, and (3) improves convergence speed. For transparency, we have included the ADAM equations as proposed by \cite{kingma2014adam} below. Default values for $\beta_1, \beta_2, \text{and } \epsilon$ used ($0.9$, $0.999$, $10^{-8}$, respectively) were used throughout. ADAM optimization was applied to each unconstrained parameter
$\theta = \{\beta, c_1, \dots, c_K, m\}$, independently at each iteration $t$. As above, the minimum is then identified iteratively with $t$ steps, learning rate $\eta$, and the derived gradients from the loss function. Optimization is performed entirely in the unconstrained, reparameterized state space to achieve a smoother optimization surface, thus increasing the stability of the model. Costs remain appropriately considered in the constrained space through back-transformations of the reparameterized values to arrive at $ \theta \in \{f_1, \dots,f_K, M\} $.


\begin{align*}
g_t &= \nabla_\theta L(\theta_t) \\
m_t &= \beta_1 m_{t-1} + (1 - \beta_1) g_t  \\
v_t &= \beta_2 v_{t-1} + (1 - \beta_2) g_t^2  \\
\hat{m}_t &= \frac{m_t}{1 - \beta_1^t}  \\
\hat{v}_t &= \frac{v_t}{1 - \beta_2^t}  \\
\theta_{t+1} &= \theta_t - \eta \frac{\hat{m}_t}{\sqrt{\hat{v}_t} + \epsilon} \\
\end{align*}



\vspace{10mm}
\bibliographystyle{unsrt}
\bibliography{ref}

\end{document}
